% ---------------------------------------------------------------------------
% Encoding, language, typography
% NOTE: sapthesis already loads xkeyval, etoolbox, geometry, fontenc, textcomp, lmodern,
% caption, graphicx, color, booktabs, amsmath, fancyhdr. Do not load them again.
% ---------------------------------------------------------------------------
\usepackage[utf8]{inputenc}
\usepackage{microtype}
\usepackage{csquotes}          % required by biblatex for proper quoting
\usepackage{tikz}
\usepackage{subcaption}

% ---------------------------------------------------------------------------
% Math, figures, tables, code
% ---------------------------------------------------------------------------
\usepackage{amssymb}                 % amsmath, graphicx, booktabs, fontenc, lmodern, caption are loaded by sapthesis
\graphicspath{{assets/plots/}{assets/icons/}}   % relative to the project root
\usepackage{tabularx}
\usepackage{placeins}        % \FloatBarrier: keep figures inside their own section
\usepackage{siunitx}           % units and number formatting: \SI{8}{\giga\byte}
\usepackage{subcaption}        % subfigures
\usepackage{xcolor}
\usepackage{tikz}
\usetikzlibrary{arrows.meta, positioning, calc, fit, backgrounds, decorations.pathreplacing}
\usepackage{pgfplots}
\pgfplotsset{compat=1.18}
\usepackage{listings}          % source excerpts from NCCL
\lstset{
  basicstyle=\ttfamily\small,
  columns=fullflexible,
  breaklines=true,
  numbers=left,
  numberstyle=\tiny,
  frame=single,
  captionpos=b,
}

% ---------------------------------------------------------------------------
% Bibliography (biblatex + biber)
% ---------------------------------------------------------------------------
\usepackage[
  backend=biber,
  style=numeric,
  sorting=nyt,
  maxbibnames=10,
]{biblatex}
\addbibresource{src/references.bib}

% ---------------------------------------------------------------------------
% Links and cross-references (load hyperref late, cleveref after it)
% ---------------------------------------------------------------------------
\usepackage[hidelinks]{hyperref}   % print-friendly; see sapthesis doc for colors
\usepackage[nameinlink,capitalise]{cleveref}

% ---------------------------------------------------------------------------
% Acronyms
% ---------------------------------------------------------------------------
\usepackage[printonlyused]{acronym}

% ---------------------------------------------------------------------------
% Project macros (keep notation consistent across chapters)
% ---------------------------------------------------------------------------
\newcommand{\bigO}[1]{\mathcal{O}\!\left(#1\right)}
\newcommand{\nccl}{\textsc{NCCL}}
\newcommand{\bine}{\textsc{Bine}}
\newcommand{\pat}{\textsc{PAT}}
\newcommand{\todo}[1]{\textcolor{red}{[TODO: #1]}}   % remove before submission


\usepackage{amsmath,amssymb,amsthm,booktabs,algorithm,algpseudocode}
\newtheorem{lemma}{Lemma}
\newcommand{\rtn}{\mathit{rank2nb}}
\newcommand{\ntr}{\mathit{nb2rank}}
\newcommand{\nbr}[1]{\overline{#1}}
\newcommand{\Send}{\mathcal{S}}
\newcommand{\Recv}{\mathcal{R}}
\newcommand{\none}{\bot}

\definecolor{crosslink}{RGB}{200,40,40}
\definecolor{localxfer}{RGB}{60,150,70}

\definecolor{chunkA}{RGB}{60,170,70}   % chunk 1 (green)
\definecolor{chunkB}{RGB}{240,200,40}  % chunk 2 (yellow)
\definecolor{chunkC}{RGB}{50,110,220}  % chunk 3 (blue)
\colorlet{idle}{gray!35}

% \ringpic{label}{color 0->1}{color 1->2}{color 2->0}
\newcommand{\ringpic}[4]{%
  \begin{tikzpicture}[baseline=(lab.north),
      rank/.style={circle, draw, minimum size=0.65cm, inner sep=0pt, fill=white},
      lnk/.style={-{Stealth[length=2.5mm]}, line width=2.2pt,
                  shorten >=1pt, shorten <=1pt}]
    \node[rank] (r0) at (90:1.1)   {0};
    \node[rank] (r1) at (-30:1.1)  {1};
    \node[rank] (r2) at (210:1.1)  {2};
    \draw[lnk, #2] (r0) -- (r1);
    \draw[lnk, #3] (r1) -- (r2);
    \draw[lnk, #4] (r2) -- (r0);
    \node (lab) at (0,-1.9) {#1};
  \end{tikzpicture}%
}

% \ftstep{holders}{transfers a/b}{uplink loads j/label}: one step of a broadcast on
% 8 ranks, 4 leaf switches and 1 spine (2:1 oversubscribed fat tree). Each transfer is
% drawn along the links it uses: green if it stays under one leaf switch, red if it
% crosses the uplinks. Ranks in {holders} already have the data; the load of each used
% uplink is written next to it.
\newcommand{\ftstep}[3]{%
  \begin{tikzpicture}[x=0.56cm, y=0.66cm,
      rank/.style={circle, draw, minimum size=0.4cm, inner sep=0pt,
                   fill=white, font=\scriptsize},
      sw/.style={rectangle, rounded corners=1.5pt, draw=black!60, fill=black!8,
                 minimum width=0.62cm, minimum height=0.3cm, inner sep=0pt},
      wire/.style={draw=black!25, line width=0.6pt},
      uplink/.style={draw=black!25, line width=1.6pt},
      path/.style={line width=1.5pt, line join=round, line cap=round},
      load/.style={font=\scriptsize\bfseries, crosslink, inner sep=1pt, fill=white}]
    \foreach \i in {0,...,7} \coordinate (r\i) at (\i,0);
    \foreach \j in {0,...,3} \coordinate (l\j) at (2*\j+0.5,1.6);
    \coordinate (sp) at (3.5,3.3);
    % idle links
    \foreach \i in {0,...,7} {\pgfmathtruncatemacro{\j}{\i/2}\draw[wire] (r\i) -- (l\j);}
    \foreach \j in {0,...,3} \draw[uplink] (l\j) -- (sp);
    % transfers, drawn along the links they use
    \foreach \a/\b in {#2} {
      \pgfmathtruncatemacro{\ja}{\a/2}
      \pgfmathtruncatemacro{\jb}{\b/2}
      \ifnum\ja=\jb
        \draw[path, localxfer, -{Stealth[length=1.8mm]}] (r\a) -- (l\ja) -- ([yshift=0.2cm]r\b);
      \else
        \draw[path, crosslink, -{Stealth[length=1.8mm]}] (r\a) -- (l\ja) -- (sp) -- (l\jb) -- ([yshift=0.2cm]r\b);
      \fi
    }
    % switches and ranks on top of the paths
    \foreach \j in {0,...,3} \node[sw] at (l\j) {};
    \node[sw, minimum width=1.2cm] at (sp) {};
    \foreach \i in {0,...,7} \node[rank] at (r\i) {\i};
    \foreach \h in {#1} \node[rank, fill=black!20] at (r\h) {\h};
    % uplink loads
    \foreach \j/\lab in {#3} \node[load] at ($(l\j)!0.42!(sp)$) {\lab};
  \end{tikzpicture}%
}

% \binepic{label}{active edges}: 8-rank distance-halving Bine tree rooted at 0.
% Tree from the XOR rule q = nb2rank(rank2nb(r) xor 1..1):
%   step 0: 0->3 | step 1: 0->7, 3->4 | step 2: 0->1, 3->2, 7->6, 4->5
\newcommand{\binepic}[2]{%
  \begin{tikzpicture}[baseline=(lab.north), x=0.85cm, y=0.75cm,
      rank/.style={circle, draw, minimum size=0.55cm, inner sep=0pt,
                   fill=white, font=\scriptsize},
      idl/.style={draw=idle, line width=1.2pt},
      act/.style={-{Stealth[length=2.4mm]}, line width=2.2pt,
                  shorten >=1pt, shorten <=1pt},
      mid/.style={-{Stealth[length=2mm]}, line width=1.2pt,
                  shorten >=1pt, shorten <=1pt},
      side/.style={line width=1.2pt, shorten >=1pt, shorten <=1pt}]
    \useasboundingbox (-1.8,0.5) rectangle (2.7,-4.5);
    \node[rank] (n0) at (0.7,0)     {0};
    \node[rank] (n3) at (-0.6,-1.1) {3};
    \node[rank] (n7) at (0.9,-1.1)  {7};
    \node[rank] (n1) at (2.1,-1.1)  {1};
    \node[rank] (n4) at (-1.3,-2.2) {4};
    \node[rank] (n2) at (-0.1,-2.2) {2};
    \node[rank] (n6) at (0.9,-2.2)  {6};
    \node[rank] (n5) at (-1.3,-3.3) {5};
    \foreach \a/\b in {0/3,0/7,0/1,3/4,3/2,7/6,4/5} \draw[idl] (n\a) -- (n\b);
    #2
    \node[font=\footnotesize, align=center] (lab) at (0.45,-4.1) {#1};
  \end{tikzpicture}%
}

% \bact{from}{to}{colour}: one chunk on a link
\newcommand{\bact}[3]{\draw[act, #3] (n#1) -- (n#2);}

% \bslice{from}{to}: the whole slice (all three chunks) on a link
\newcommand{\bslice}[2]{%
  \draw[side, chunkA, transform canvas={xshift=-2.4pt}] (n#1) -- (n#2);
  \draw[side, chunkC, transform canvas={xshift=2.4pt}]  (n#1) -- (n#2);
  \draw[mid, chunkB] (n#1) -- (n#2);}

% \bcastrow{chunk letter}{colour}{first time slot}: one row of the broadcast figure
\newcommand{\bcastrow}[3]{%
  \binepic{chunk #1, step 0\\$t{=}#3$}{\bact{0}{3}{#2}}\hfill
  \binepic{chunk #1, step 1\\$t{=}\the\numexpr#3+1\relax$}%
    {\bact{0}{7}{#2}\bact{3}{4}{#2}}\hfill
  \binepic{chunk #1, step 2\\$t{=}\the\numexpr#3+2\relax$}%
    {\bact{0}{1}{#2}\bact{3}{2}{#2}\bact{7}{6}{#2}\bact{4}{5}{#2}}}

% \blab{from}{to}{slot}{colour}: link used at a given time slot
\newcommand{\blab}[4]{%
  \draw[act, #4] (n#1) -- node[pos=0.55, fill=white, inner sep=0.6pt,
      font=\tiny\bfseries] {#3} (n#2);}

% \bchunk{letter}{slots before this chunk}{colour}: whole tree of one chunk
\newcommand{\bchunk}[3]{%
  \binepic{chunk #1\\slots \the\numexpr#2+1\relax--\the\numexpr#2+3\relax}{%
    \blab{0}{3}{\the\numexpr#2+1\relax}{#3}%
    \blab{0}{7}{\the\numexpr#2+2\relax}{#3}%
    \blab{3}{4}{\the\numexpr#2+2\relax}{#3}%
    \blab{0}{1}{\the\numexpr#2+3\relax}{#3}%
    \blab{3}{2}{\the\numexpr#2+3\relax}{#3}%
    \blab{7}{6}{\the\numexpr#2+3\relax}{#3}%
    \blab{4}{5}{\the\numexpr#2+3\relax}{#3}}}

% Preamble (colors; the tikz libraries arrows.meta, positioning, fit, backgrounds, calc are already loaded)
\definecolor{ch0}{RGB}{60,130,210}
\definecolor{ch1}{RGB}{225,120,50}
\definecolor{ch2}{HTML}{27AE60}
